% ---------------------------------------------------------------------
% Front matter. Two versions controlled by \ifanonymous in main.tex.



\begin{titlepage}
\centering
\vspace*{2cm}
{\LARGE\bfseries \thetitle \par}
\vspace{2cm}

\ifanonymous
  {\large Exam number: \examnumber \par}
  \vspace{0.5em}
  {\large MSc Economics Dissertation \par}
  \vspace{0.5em}
  {\large \today \par}
  \vspace{2cm}
  {\normalsize Word count: \wordcount \par}
  \vspace{1em}
  {\normalsize Data files used: O*NET 30.3 (public); Annual Population Survey 2022--2025 (public); code and scoring archive submitted separately to \texttt{sgpe@ed.ac.uk}. \par}
\else
  {\large Tommy Ranyard \par}
  \vspace{0.5em}
  {\large MSc Economics, University of Edinburgh \par}
  \vspace{0.5em}
  {\large \today \par}
  \vspace{2cm}
  \begin{minipage}{0.85\textwidth}
  \small
  This dissertation is undertaken in collaboration with the Scottish Fiscal Commission (SFC) through the SGPE Sponsored Dissertation Programme. I thank my academic supervisors, Mark Schaffer and Eoin McLaughlin, along with Silvia Palombi and Alexander Sibelle at the SFC, for their helpful comments, and the OBR for insight into their approach, prompt, and results. All errors are my own.
  \end{minipage}
\fi

\vfill
\end{titlepage}

% ---------------------------------------------------------------------
\begin{abstract}
\noindent
Artificial intelligence is expected to reshape labour demand unevenly across occupations, and therefore across regions whose occupational composition differs. Whether Scotland is differentially exposed relative to the rest of the UK (rUK) is a fiscal question, yet exposure measures report point estimates without characterising their dependence on the scoring model. This dissertation makes three contributions. First, I prove that the regional differential in an employment-weighted exposure index removes the common component of any model-specific bias, that an affine bias survives only as a rescaling which a standardised or rank-based differential removes, and that the remainder is Cauchy--Schwarz bounded. Second, I answer it, and test the answer by re-scoring all 17,945 O*NET tasks on two further models. Scotland sits about two percentile points below the rUK in the UK occupational exposure distribution, or 0.86 index points on a 0--100 scale (95\% interval $[-1.07, -0.69]$), with London accounting for two-fifths of that gap and Scottish exposure tilting towards substitution. The sign and percentile magnitude hold across models where the index-point figure does not, the alternatives compressing the scale rather than reordering occupations. Third, the gap implies a ten-year TFP differential under three basis points and a net-position sensitivity of about half a million pounds a year. The lesson is that scale-free differentials can be trusted where levels cannot.

\medskip
\noindent\textbf{Keywords:} artificial intelligence, labour-market exposure, occupational composition, Scotland, shift-share, measurement error.\\
\textbf{JEL classification:} J21, J24, O33, R23.
\end{abstract}
\clearpage